%% Report-privacy analysis moved from the main text (Section 3.2).
Each client $k$ releases its two group means with Gaussian noise, $\tilde{\mu}_{k,a} = \mu_{k,a} + \xi_{k,a}$ with $\xi_{k,a} \sim \mathcal{N}(0, \sigma_k^2)$, calibrated by R\'enyi accounting as described in the main text, and reports $\widetilde{\dpd}_k = |\tilde{\mu}_{k,0} - \tilde{\mu}_{k,1}|$.

\begin{lemma}[Upward bias of privatised disparity]
\label{lem:folded_normal}
Let $\tilde{\mu}_{k,a} \sim \mathcal{N}(\mu_{k,a}, \sigma_k^2)$ independently. Under demographic parity ($\mu_{k,0} = \mu_{k,1}$), $\widetilde{\dpd}_k$ follows a half-normal distribution~\citep{leone1961folded} with
\begin{equation}
\label{eq:folded_exp}
\EE[\widetilde{\dpd}_k] = \tilde{\sigma}_k \sqrt{\frac{2}{\pi}} > 0, \qquad \tilde{\sigma}_k = \sqrt{2}\,\sigma_k .
\end{equation}
For weights $\bm{w} \in \Delta_K$ that are fixed or determined by earlier rounds (measurable with respect to the filtration $\mathcal{F}_{t-1}$),
\begin{equation}
\label{eq:accum_bias}
\EE\left[ \sum_{k=1}^K w_k \widetilde{\dpd}_k \;\middle|\; \mathcal{F}_{t-1} \right] = \sqrt{\frac{2}{\pi}} \sum_{k=1}^K w_k \tilde{\sigma}_k > 0,
\end{equation}
and for weights chosen after seeing the current reports, $\sum_k w_k \widetilde{\dpd}_k \ge \min_k \widetilde{\dpd}_k$, whose expectation is also positive. The noise therefore accumulates across clients instead of cancelling.
\end{lemma}
\begin{proof}
$Z = \tilde{\mu}_{k,0} - \tilde{\mu}_{k,1} \sim \mathcal{N}(0, 2\sigma_k^2)$, so $|Z|$ is half-normal with scale $\tilde{\sigma}_k$ and mean $\tilde{\sigma}_k\sqrt{2/\pi}$. Taking the conditional expectation with $\mathcal{F}_{t-1}$-measurable weights gives Eq.~\eqref{eq:accum_bias}. For arbitrary weights on the simplex, the weighted sum of non-negative reports is at least their minimum, which has positive mean.
\end{proof}

The bias comes from the absolute value. A client could instead release the signed difference $\tilde{\mu}_{k,0} - \tilde{\mu}_{k,1}$, which is unbiased. The next result concerns a narrow, single-round Gaussian testing problem and is not a general impossibility result for private fairness mechanisms.

\begin{proposition}[Per-round observability bound under local differential privacy]
\label{thm:minimax_lecam}
Let the clients release privatised disparity reports $\widetilde{\bm{D}} = (\widetilde{\dpd}_1, \dots, \widetilde{\dpd}_K)$ generated by this release with effective noise standard deviation $\tilde{\sigma}$. For any decision rule $\psi(\widetilde{\bm{D}}) \in \{0, 1\}$ that uses one round of reports to distinguish $H_0$: every client's disparity satisfies $\delta_k \le \tau/2$, from $H_1$: every $\delta_k \ge 3\tau/2$, the minimax testing error $\Rstar = \inf_\psi \max\{ \Pr_{H_0}[\psi=1], \Pr_{H_1}[\psi=0] \}$ satisfies
\begin{equation}
\label{eq:lecam_bound}
\Rstar \;\ge\; \frac{1}{2} \left( 1 - \frac{\tau \sqrt{K}}{2 \tilde{\sigma}} \right).
\end{equation}
\end{proposition}
\begin{proof}
Take the two hypotheses $\delta_k^{(0)} = \tau/2$ and $\delta_k^{(1)} = 3\tau/2$ for every client. Le Cam's two-point lemma~\citep{lecam1973convergence, tsybakov2009nonparametric} gives
$\inf_\psi \max\{ \Pr_{H_0}[\psi=1], \Pr_{H_1}[\psi=0] \} \ge \frac{1}{2}( 1 - \mathrm{TV}(P_0, P_1))$.
Folding is a deterministic map, so by the data processing inequality $\mathrm{TV}(P_0, P_1) \le \mathrm{TV}(P_0^{\mathrm{Gauss}}, P_1^{\mathrm{Gauss}})$. For $K$ independent Gaussians with mean separation $\tau$ and variance $\tilde{\sigma}^2$,
$\mathrm{KL}(P_0^{\mathrm{Gauss}} \,\|\, P_1^{\mathrm{Gauss}}) = K\tau^2/(2\tilde{\sigma}^2)$, and Pinsker's inequality $\mathrm{TV} \le \sqrt{\mathrm{KL}/2}$ gives $\mathrm{TV} \le \tau\sqrt{K}/(2\tilde{\sigma})$.
\end{proof}

\begin{remark}[Scope of the bound]
\label{rem:lecam_scope}
The bound concerns one round, a fixed two-point hypothesis, independent Gaussian releases, and a common effective noise level. Pooling repeated rounds or modeling heterogeneous client noise can materially change it. Table~\ref{tab:c21_privacy_diagnostic} therefore reports the heterogeneous calculation as a diagnostic approximation. The model-update channel is outside this analysis.
\end{remark}
